% day1_seg4_prob/prob_measures.tex
\begin{frame}
  \onslide<1->{
	\frametitle{Probability} 
	\framesubtitle{Set Functions and Measures}
}
\begin{itemize}
  \onslide<2->{
\item A measure is a way of assigning a number to a set.  
}
  \onslide<3->{
\item Intuitively, it measures the size of the set.  
}
  \onslide<4->{
\item In probability, a measure is a way of assigning a probability to an event.  
}
\end{itemize}
\end{frame}

\begin{frame}
  \onslide<1->{
	\frametitle{Probability} 
	\framesubtitle{Set Functions and Measures}
}
\begin{itemize}
  \onslide<2->{
\item  A set function $\mu$ on $\mathcal{M}$ assigns a real number $\mu(A)$ to each set $A\in\mathcal{M}$. 
}
  \onslide<3->{
\item A set function $\mu$ on $\mathcal{M}$ is said to be \textbf{\textit{additive}}, if $\mu(A\cup B)=\mu(A)+\mu(B)$ whenever $A$ and $B$ are disjoint.  
}
  \onslide<4->{
\item $\mu$ is called \textbf{\textit{countably additive}} if $\mu(\bigcup_{i\in I}A_i)=\sum_{i\in I}\mu(A_i)$, if the family of sets $\set{A_i:i\in I}$ is disjoint.  
}
\end{itemize}
\end{frame}

\begin{frame}
  \onslide<1->{
	\frametitle{Probability} 
	\framesubtitle{Set Functions and Measures}
}
\begin{itemize}
  \onslide<2->{
\item A measure $\mu$ on a $\sigma$-algebra $\mathcal{M}$ is a real-valued, non-negative, countably additive set function:  $\mu:\mathcal{M}\rightarrow\Real$  
}
  \onslide<3->{
\item Measures are used in mathematics to define abstract integration.  
}
\end{itemize}
\end{frame}

\begin{frame}
  \onslide<1->{
	\frametitle{Probability} 
	\framesubtitle{Set Functions and Measures}
}
\begin{itemize}
  \onslide<2->{
\item Example: Counting Measure. $\Omega$ is finite and the measure of each subset is merely the number of elements in the set.  
}
  \onslide<3->{
\item Example: Dirac Measure.  Let $\mathcal{M}$ be any $\sigma$-algebra and
let $\omega\in \Omega$ be an arbitrary but fixed element of
$\Omega$. For each $A\in\mathcal{M}$ define
  \[ \mu_{\omega}(A) = \left\lbrace
  \begin{array}{c l}
    1 & \text{if $\omega\in A$},\\
    0 & \text{otherwise}.
  \end{array}
\right. \]  
}
\end{itemize}
\end{frame}

\begin{frame}
  \onslide<1->{
	\frametitle{Probability} 
	\framesubtitle{Probability Measures}
}
\begin{itemize}
  \onslide<2->{
\item A measure $\mu$ on $\mathcal{M}$ with the property 
 that $\mu(\Omega)=1$ is called a \textbf{probability measure} or \textbf{probability function}.  
}
  \onslide<3->{
\item Instead of $\mu$ we will use $P$, which is
 more customary.  
}
\end{itemize}
\end{frame}

\begin{frame}
  \onslide<1->{
	\frametitle{Probability} 
	\framesubtitle{Probability: Specialized Terminology}
}
\begin{itemize}
  \onslide<2->{
\item Null Event:  $P(A)=0$.  It is important to realize this
 does not imply that $A=\emptyset$.  In measure theory this would
 be called a set of measure zero.  
}
  \onslide<3->{
\item Independent Events:  $P(A\cap B)=P(A)P(B)$  
}
  \onslide<4->{
\item Mutually Exclusive Events: \\
 $A\cap B=\emptyset\Rightarrow P(A\cup B)=P(A)+P(B)$  
}
\end{itemize}
\end{frame}

\begin{frame}
  \onslide<1->{
	\frametitle{Probability} 
	\framesubtitle{Examples: Probability Measure}
}
\begin{itemize}
  \onslide<2->{
\item Double coin toss:  
\begin{itemize}
}  
  \onslide<3->{
\item Sample Space: $\Omega=\set{HH,HT,TH,TT}$  
}
  \onslide<4->{
\item $\sigma$-algebra: $\mathcal{M}=\set{\emptyset, \set{HH},\set{HT,TH,TT},\Omega}$  
}
  \onslide<5->{
\item Probability Measure: $P(A)=$number of elements in $A$ divided by total number of elements in $\Omega$
}
  \onslide<6->{
\item E.g., $P(\set{HH})=1/4$
}
\end{itemize}
\end{itemize}
\end{frame}
\begin{frame}
  \onslide<1->{
	\frametitle{Probability} 
	\framesubtitle{Examples: Probability Measure}
}
\begin{itemize}
  \onslide<2->{
\item Let $\Omega$ be any finite set of cardinality $n$.  Let the
 $\sigma$-algebra be the set of all subsets of $\Omega$ and let
 $|\cdot|$ denote the counting measure, i.e. for $A\subseteq\Omega$,
 $\|A|=$number of elements in $A$.  Define a probability measure
 by $P=\frac{1}{n}|\cdot|$.  
}
  \onslide<3->{
\item To illustrate the last example a little more, let
 $\Omega=\set{\omega_1,\ldots,\omega_6}$, and let
 $A=\set{\omega_1,\omega_2}$.  Then $P(A)=\frac{1}{6}\cdot|A|=\frac{1}{6}\cdot 2=\frac{1}{3}$  
}
\end{itemize}
\end{frame}
\begin{frame}
  \onslide<1->{
	\frametitle{Probability} 
	\framesubtitle{Examples: Probability Measure}
}
\begin{itemize}
  \onslide<2->{
\item Let $\Omega=[0,1]$ and let $\mathcal{B}$ be the Borel
 $\sigma$-algebra on the interval $[0,1]$.  Define $P=\lambda$,
 the Lebesgue measure on $\mathcal{B}$.  Note that for any
 open subinterval $(a,b)$ of $[a,b]$, we have $P(a,b)=b-a$.  
}
\end{itemize}
\end{frame}
\begin{frame}
  \onslide<1->{
	\frametitle{Probability} 
	\framesubtitle{Conditional Probability}
}
\begin{itemize}
  \onslide<2->{
\item Suppose you have two events $A$ and $B$ and you know that the event $B$ has in fact occured.    
}
  \onslide<3->{
\item How does that knowledge affect the probability that both $A$ and $B$ have occured?  
}
  \onslide<4->{
\item Intuitively, we would tend to think that the likelihood of $A\cap B$ increases with the knowledge that $B$ has occured.  
}
\end{itemize}
\end{frame}


\begin{frame}
  \onslide<1->{
	\frametitle{Probability} 
	\framesubtitle{Conditional Probability}
}
\begin{itemize}
  \onslide<2->{
\item The increased likelihood of $A$ given 
that $B$ has occured is called the \textbf{conditional probability} of $A$ given $B$.   
}
  \onslide<3->{
\item If $P(B)\neq 0$, then the conditional probability of $A$,
 given that $B$ has occurred, is:  $P(A|B):=P(A\cap B)/P(B)$  
}
  \onslide<4->{
\item Sometimes we say the probability of $A$ conditioned on $B$.  
}
  \onslide<5->{
\item Note that since $P(B)\leq 1$, $P(A\cap B)/P(B)\geq P(A\cap B)$  
}
\end{itemize}
\end{frame}

\begin{frame}
  \onslide<1->{
	\frametitle{Probability} 
	\framesubtitle{Conditional Probability: Bayes' Theorem}
}
\begin{itemize}
  \onslide<2->{
\item Since $P(A|B)=P(A\cap B)/P(B)$ and $P(B|A):=P(B\cap A)/P(A)$, we can eliminate $P(A\cap B)$ and derive
\[P(A|B)=\dfrac{P(B|A)P(A)}{P(B)}\]  
}
  \onslide<3->{
\item Using this we obtain Bayes' Theorem (simple form):
\begin{equation}\label{bayes}
   P(A|B)=\frac{P(A)P(B|A)}{P(A)P(B|A)+P(A')P(B|A')}.
\end{equation}  
}
\end{itemize}
\end{frame}

\begin{frame}
  \onslide<1->{
	\frametitle{Probability} 
	\framesubtitle{Conditional Probability}
}
\begin{itemize}
  \onslide<2->{
\item Bayes' Theorem (complex form):  Let $A_1,\ldots,A_n$ be a
 partition of $\Omega$.  Then,
\begin{equation}\label{bayescomplex}
   P(A_i|B)=\frac{P(A_i)P(B|A_i)}{\sum_{j=1}^{n}P(A_j)P(B|A_j)}.
\end{equation}  
}
  \onslide<3->{
\item Bayes' Theorem expresses posterior probability in terms of
 prior probabilities and conditional probabilities.  
}
\end{itemize}
\end{frame}
\begin{frame}
  \onslide<1->{
	\frametitle{Probability} 
	\framesubtitle{Conditional Probability}
}
\begin{itemize}
  \onslide<2->{
\item In one derivation of the Kalman filter, the use of Bayes'
 Theorem is a crucial step.  We'll see this later.  
}
  \onslide<3->{
\item Bayes' Theorem is used to derive the Kalman gain: to
 determine the best estimate of state \textbf{after} measurement has
 occurred.  
}
\end{itemize}
\end{frame}

\begin{frame}
  \onslide<1->{
	\frametitle{Probability} 
	\framesubtitle{Conditional Probability}
}
\begin{itemize}
  \onslide<2->{
\item Bayes' Theorem allows us to reverse the conditional probability.  
}
  \onslide<3->{
\item We can derive information about the state conditioned on the measurement from 
information about the measurement conditioned on the state.    
}
\end{itemize}
\end{frame}
